\documentclass[11pt,a4paper]{report}

% ============================================================
% PAKETE
% ============================================================
\usepackage[utf8]{inputenc}
\usepackage[T1]{fontenc}
\usepackage[ngerman]{babel}
\usepackage{amsmath, amssymb}
\usepackage{geometry}
\usepackage{graphicx}
\usepackage{booktabs}
\usepackage{array}
\usepackage{longtable}
\usepackage{hyperref}
\usepackage{xcolor}
\usepackage{listings}
\usepackage{enumitem}
\usepackage{tcolorbox}
\usepackage{tikz}
\usepackage{float}
\usepackage{microtype}
\usepackage{lmodern}
\usepackage{setspace}

% ============================================================
% LAYOUT
% ============================================================
\geometry{
	a4paper,
	top=2.5cm, bottom=2.5cm,
	left=2.5cm, right=2.5cm
}
\setlength{\parskip}{6pt}
\setlength{\parindent}{0pt}
\onehalfspacing

% ============================================================
% FARBEN & BOXEN
% ============================================================
\definecolor{boxblue}{RGB}{220,235,250}
\definecolor{boxgreen}{RGB}{220,245,220}
\definecolor{codegray}{RGB}{245,245,245}
\definecolor{codegreen}{rgb}{0,0.5,0}
\definecolor{codepurple}{rgb}{0.5,0,0.5}

\tcbuselibrary{skins, breakable}

\newtcolorbox{analogiebox}[1][]{
	colback=boxblue, colframe=blue!50!black,
	title=Analogie, fonttitle=\bfseries,
	breakable, #1
}

\newtcolorbox{merkbox}[1][]{
	colback=boxgreen, colframe=green!40!black,
	title=Merksatz, fonttitle=\bfseries,
	breakable, #1
}

\newtcolorbox{hinweisbox}[1][]{
	colback=yellow!15, colframe=orange!70!black,
	title=Hinweis, fonttitle=\bfseries,
	breakable, #1
}

% ============================================================
% LISTINGS (Code)
% ============================================================
\lstset{
	backgroundcolor=\color{codegray},
	basicstyle=\ttfamily\small,
	breaklines=true,
	captionpos=b,
	commentstyle=\color{codegreen},
	keywordstyle=\color{blue}\bfseries,
	stringstyle=\color{codepurple},
	numbers=left,
	numberstyle=\tiny\color{gray},
	frame=single,
	rulecolor=\color{gray!40},
	tabsize=4,
	showstringspaces=false,
	literate={ä}{{\"{a}}}1 {ö}{{\"{o}}}1 {ü}{{\"{u}}}1
	{Ä}{{\"{A}}}1 {Ö}{{\"{O}}}1 {Ü}{{\"{U}}}1
	{ß}{{\ss{}}}1
}

% ============================================================
% HYPERREF
% ============================================================
\hypersetup{
	colorlinks=true,
	linkcolor=blue!60!black,
	urlcolor=blue!60!black,
	citecolor=blue!60!black,
	pdftitle={HILOMOT -- Methode und Tool-Dokumentation},
	pdfauthor={heinzl75},
	pdfsubject={Lokale Modellnetze, HILOMOT, LOLIMOT}
}

% ============================================================
% DOKUMENT
% ============================================================


\title{\textbf{HILOMOT -- Einführung, Motivation und Toolbeschreibung}\\\large}
\author{}
\date{}

\begin{document}
	\maketitle
	\tableofcontents
	\newpage
	
	
	% ============================================================
	\section{Einführung und Motivation}
	% ============================================================
	
	\textbf{HILOMOT} (\textit{Hierarchical Local Model Tree}) ist ein Algorithmus zum Aufbau lokaler Modellnetze (LMN). Er partitioniert den Eingangsraum schrittweise durch binäre Splits und passt in jeder Region ein lineares Modell an.
	
	In der Praxis begegnen uns ständig \textbf{nichtlineare Systeme} -- ob Verbrennungsmotoren, chemische Reaktoren oder Roboterkinematiken. Ein einfaches lineares Modell $y = a \cdot x + b$ kann solche Zusammenhänge nicht ausreichend beschreiben.
	
	\textbf{Grundidee:} Statt ein einziges kompliziertes nichtlineares Modell zu bauen, zerlegen wir das Problem in \textbf{viele kleine Teilprobleme}, die jeweils einfach (linear) lösbar sind.
	
	\begin{analogiebox}
		Stell dir eine Landkarte vor. Statt die gesamte Erdoberfläche mit einer einzigen Formel zu beschreiben, verwenden wir viele kleine, flache Kartenbla"tter. Jedes Blatt ist für sich genommen einfach (flach/linear), aber zusammen beschreiben sie die gekrümmte Oberfläche beliebig genau.
	\end{analogiebox}
	
	Im Gegensatz zum Vorgänger \textbf{LOLIMOT} verwendet HILOMOT \textbf{oblique} (schräge) Trennhyperebenen, die nicht auf die Koordinatenachsen beschränkt sind. Damit kann das Modell auch diagonale und nichtachsenparallele Strukturen in den Daten effizient erfassen.
	
	Dieses Python-Tool ist eine Portierung der MATLAB LMNTool-Toolbox (Version 1.5.2.1) und stellt sowohl eine interaktive GUI als auch eine skriptbasierte Schnittstelle bereit.
	
	\begin{figure}[H]
		\centering
		\includegraphics[width=0.85\textwidth]{hilomot_3d_partition_comparison.png}
		\caption{Typische Ausgabe des Python-Tools mit modellierter 3D-Ergebnisfläche und zugehörig em Partitionierungsplot.}
		\label{fig:intro}
	\end{figure}
	
	% ============================================================
	\section{Theoretische Grundlagen}
	% ============================================================
	
	\subsection{Ausgangsproblem: Warum lokale Modelle?}
	
	Die zentrale Frage beim Aufbau eines lokalen Modellnetzes ist: \textbf{Wie teilen wir den Eingangsraum sinnvoll auf?}
	
	Das führt uns zum Konzept der \textbf{lokalen Modellnetze:}
	
	\begin{equation}
		\hat{y}(x) = \sum_{i=1}^{M} \Phi_i(x) \cdot \hat{y}_i(x)
	\end{equation}
	
	\begin{itemize}
		\item $\hat{y}_i(x)$: ein einfaches lineares Modell, das in einer bestimmten Region \glqq zuständig\grqq{} ist
		\item $\Phi_i(x)$: eine Gewichtungsfunktion, die angibt, \textbf{wie stark} Modell $i$ an der Stelle $x$ beiträgt
		\item Es gilt immer: $\sum_i \Phi_i(x) = 1$ -- die Gewichte summieren sich zu~1 (\textit{Partition of Unity})
	\end{itemize}
	
	% ----------------------------------------
	\subsection{Lokale Modellnetze (LMN)}
	% ----------------------------------------
	
	Ein lokales Modellnetz approximiert eine nichtlineare Funktion $f: \mathbb{R}^n \rightarrow \mathbb{R}$ durch eine gewichtete Summe lokaler Modelle:
	
	\begin{equation}
		\hat{y}(\mathbf{x}) = \sum_{i=1}^{M} \Phi_i(\mathbf{z}) \cdot \hat{y}_i(\mathbf{x})
	\end{equation}
	
	wobei:
	\begin{itemize}
		\item $M$ die Anzahl der lokalen Modelle (Blätter) ist,
		\item $\hat{y}_i(\mathbf{x}) = \boldsymbol{\theta}_i^T \tilde{\mathbf{x}}$ das $i$-te lokale lineare Modell,
		\item $\tilde{\mathbf{x}} = [1, x_1, \ldots, x_n]^T$ der erweiterte Regressionsvektor,
		\item $\Phi_i(\mathbf{z})$ die Gültigkeitsfunktion des $i$-ten Modells ist.
	\end{itemize}
	
	Die Gültigkeitsfunktionen erfüllen zu jedem Zeitpunkt die \textbf{Partition-of-Unity-Bedingung}:
	
	\begin{equation}
		\sum_{i=1}^{M} \Phi_i(\mathbf{z}) = 1 \quad \text{für alle } \mathbf{z}
	\end{equation}
	
	Der Eingangsvektor der Gültigkeitsfunktionen $\mathbf{z}$ kann identisch mit $\mathbf{x}$ sein oder davon abweichen (z.\,B. bei zeitverzögerten Systemen).
	
	% ----------------------------------------
	\subsection{Gültigkeitsfunktionen und Sigmoid-Splits}
	% ----------------------------------------
	
	Die Gültigkeitsfunktionen werden durch einen \textbf{binären Baum} aufgebaut. An jedem inneren Knoten des Baums wird der Raum durch eine Sigmoid-Funktion aufgeteilt:
	
	\begin{equation}
		\Psi(\mathbf{z}) = \sigma\left(-\frac{\kappa}{\rho} \cdot (\mathbf{w}^T \mathbf{z} + b)\right),
		\qquad \sigma(x) = \frac{1}{1+e^{-x}}
	\end{equation}
	
	wobei:
	\begin{itemize}
		\item $\mathbf{w}$ der Normalenvektor der Trennhyperebene,
		\item $b$ der Bias (Verschiebungsterm),
		\item $\kappa$ der \textbf{Schärfeparameter} (kontrolliert die Steilheit der Sigmoid-Funktion),
		\item $\rho$ der \textbf{Glättigkeitsparameter} (\texttt{smoothness}, kontrolliert die überlappung der Regionen).
	\end{itemize}
	
	Diese Funktion gibt einen Wert zwischen 0 und~1 aus:
	\begin{itemize}
		\item Nahe~1 $\Rightarrow$ der Punkt gehört \glqq eher links\grqq{}
		\item Nahe~0 $\Rightarrow$ der Punkt gehört \glqq eher rechts\grqq{}
		\item Genau~0{,}5 $\Rightarrow$ der Punkt liegt auf der Trennfläche
	\end{itemize}
	
	\begin{hinweisbox}
		Der übergang ist \textbf{weich/glatt}, nicht hart. Es gibt keinen abrupten Sprung an der Grenze -- das sorgt für ein glattes Gesamtmodell.
	\end{hinweisbox}
	
	Die Gültigkeit eines Blattknotens ergibt sich durch das Produkt der Sigmoid-Werte entlang des Pfades von der Wurzel zum Blatt:
	
	\begin{equation}
		\Phi_i(\mathbf{z}) = \prod_{k \in \mathrm{Pfad}(i)} \Psi_k(\mathbf{z})
		\quad \text{(linkes Kind)}
	\end{equation}
	\begin{equation}
		\Phi_i(\mathbf{z}) = \prod_{k \in \mathrm{Pfad}(i)} \bigl(1 - \Psi_k(\mathbf{z})\bigr)
		\quad \text{(rechtes Kind)}
	\end{equation}
	
	% ----------------------------------------
	\subsection{LOLIMOT -- Achsenparallele Splits}
	% ----------------------------------------
	
	\textbf{LOLIMOT} (\textit{Local Linear Model Tree}) ist der Vorgänger von HILOMOT und beantwortet die Partitionierungsfrage mit einer einfachen Strategie: \textbf{Schneide den Raum immer entlang einer Achse durch.}
	
	Jeder Split erfolgt achsenparallel:
	
	\begin{equation}
		\mathbf{w} = \mathbf{e}_d, \qquad b = -c_d
	\end{equation}
	
	wobei $\mathbf{e}_d$ der $d$-te Einheitsvektor ist. Dies entspricht einer Aufteilung entlang einer Koordinatenachse bei Schwellenwert~$c_d$.
	
	\textbf{Vorgehen (vereinfacht):}
	\begin{enumerate}
		\item Starte mit einem einzigen Modell für den ganzen Raum
		\item Finde die Region mit dem grö\ss{}ten Fehler
		\item Teile sie in der Mitte -- probiere jede Achse durch, nimm den besten Schnitt
		\item Trainiere die lokalen Modelle neu
		\item Wiederhole, bis der Fehler klein genug ist
	\end{enumerate}
	
	\textbf{Stärken:} Einfach, schnell, robust, gut verstanden.
	
	\begin{hinweisbox}
		\textbf{Aber:} Was passiert, wenn die \glqq wahre\grqq{} Grenze zwischen zwei Verhaltensbereichen \textbf{schräg} verläuft?
		
		Beispiel: Ein Motor verhält sich unterschiedlich je nach Kombination von Drehzahl und Last. Die Grenze zwischen \glqq magerem\grqq{} und \glqq fettem\grqq{} Betrieb verläuft diagonal im Drehzahl-Last-Kennfeld. LOLIMOT muss diese diagonale Grenze durch viele kleine Treppenstufen annähern -- das kostet \textbf{unnötig viele Teilmodelle}.
	\end{hinweisbox}
	
	% ----------------------------------------
	\subsection{HILOMOT -- Oblique Splits}
	% ----------------------------------------
	
	HILOMOT erweitert LOLIMOT durch \textbf{oblique} (schräge) Trennhyperebenen. Statt nur achsenparallel zu schneiden, darf die Trennlinie \textbf{beliebig orientiert} sein:
	
	\begin{equation}
		\mathbf{w}^T \mathbf{z} + b = 0 \qquad \text{(allgemeine Hyperebene)}
	\end{equation}
	
	\begin{center}
		\begin{tabular}{lll}
			\toprule
			& \textbf{LOLIMOT} & \textbf{HILOMOT} \\  \midrule
			Schnittrichtung & Nur entlang der Achsen & Beliebig orientiert (schräg) \\
			Trennfläche & $x_j = \mathrm{const}$ & $\mathbf{w}^\top \mathbf{x} + b = 0$ \\
			\bottomrule
		\end{tabular}
	\end{center}
	
	\begin{analogiebox}
		LOLIMOT schneidet den Kuchen nur waagrecht oder senkrecht. HILOMOT darf auch diagonal schneiden -- und trifft damit die natürliche Form des Kuchens oft viel besser.
	\end{analogiebox}
	
	\textbf{Split-Kandidaten} werden in dieser Implementierung generiert durch:
	\begin{enumerate}
		\item Alle achsenparallelen Richtungen $\mathbf{e}_1, \ldots, \mathbf{e}_n$ (wie LOLIMOT),
		\item Falls \texttt{oblique=True}: zusätzlich die Hauptachsenrichtung der gewichteten Kovarianzmatrix der Daten im schlechtesten Blatt (PCA-Richtung).
	\end{enumerate}
	
	Zusätzlich wird bei 2D-Eingaben eine \textbf{kontinuierliche Optimierung} der Split-Richtung durchgeführt: Der Winkel~$\theta$ wird so gewählt, dass der globale Trainingsfehler minimiert wird:
	
	\begin{equation}
		\theta^* = \arg\min_\theta \mathrm{MSE}\left(\mathbf{y},\, \hat{\mathbf{y}}(\theta)\right),
		\qquad \mathbf{w}(\theta) = [\cos\theta,\, \sin\theta]^T
	\end{equation}
	
	\begin{merkbox}
		HILOMOT enthält LOLIMOT als Spezialfall. Wenn man die Trennrichtung $\mathbf{w}$ auf Einheitsvektoren $\mathbf{e}_j$ einschränkt, wird jeder HILOMOT-Split zu einem LOLIMOT-Split.
	\end{merkbox}
	
	% ----------------------------------------
	\subsection{Der Baum als Organisationsstruktur}
	% ----------------------------------------
	
	HILOMOT organisiert die Aufteilung als \textbf{Binärbaum:}
	
	\begin{lstlisting}[caption={Beispielbaum mit zwei Gate-Knoten und drei Blättern},language={}
		]
		[Wurzel-Gate]
		/             \\
		[Gate]              Modell 3
		/    \\
		Modell 1   Modell 2
	\end{lstlisting}
	
	\begin{itemize}
		\item \textbf{Innere Knoten} = Entscheidungen (Gates): \glqq Gehört der Punkt eher links oder rechts?\grqq{}
		\item \textbf{Blätter} = Lokale lineare Modelle
	\end{itemize}
	
	An jedem inneren Knoten sitzt eine \textbf{weiche Entscheidungsfunktion} (Sigmoid):
	
	\begin{equation}
		g(\mathbf{x}) = \frac{1}{1 + e^{-(\mathbf{w}^T \mathbf{x} + b)}}
	\end{equation}
	
	Das Gewicht eines Blattes ergibt sich aus dem \textbf{Produkt aller Entscheidungen auf dem Weg von der Wurzel zum Blatt.} Für Modell~1 im Beispielbaum oben:
	\begin{itemize}
		\item An der Wurzel: \glqq nach links\grqq{} $\Rightarrow$ Gewicht $g_{\mathrm{Wurzel}}(\mathbf{x})$
		\item Am zweiten Knoten: \glqq nach links\grqq{} $\Rightarrow$ Gewicht $g_{\mathrm{Knoten}}(\mathbf{x})$
		\item Gesamt: $\Phi_1(\mathbf{x}) = g_{\mathrm{Wurzel}}(\mathbf{x}) \cdot g_{\mathrm{Knoten}}(\mathbf{x})$
	\end{itemize}
	
	Dieses Pfadprodukt garantiert automatisch, dass sich alle Blattgewichte zu~1 summieren.
	
	% ----------------------------------------
	\subsection{Kappa-Anpassung (\texttt{demanded\_membership\_value})}
	% ----------------------------------------
	
	Der Schärfeparameter $\kappa$ beeinflusst, wie stark sich die Regionen überlappen. Ein hoher Wert erzeugt scharfe, nahezu unstetige übergänge; ein niedriger Wert erzeugt weiche, breite übergänge.
	
	Die Kappa-Anpassung basiert auf dem \textbf{demanded membership value} $\mu^*$ (Standard: 0{,}8). $\kappa$ wird so gewählt, dass Datenpunkte in der Mitte ihrer Region mindestens die Gültigkeitsfunktion~$\mu^*$ erreichen:
	
	\begin{equation}
		\kappa = \left|\frac{\rho \cdot \ln\left(\dfrac{1}{\mu^*} - 1\right)}{\mathbf{w}^T \mathbf{z}_{\min} + b}\right|
	\end{equation}
	
	wobei $\mathbf{z}_{\min}$ derjenige Datenpunkt ist, der der Trennhyperebene am nächsten liegt.
	
	% ----------------------------------------
	\subsection{Gewichtete kleinste Quadrate (WLS)}
	% ----------------------------------------
	
	Die Parameter $\boldsymbol{\theta}_i$ jedes lokalen Modells werden durch \textbf{gewichtete kleinste Quadrate} (WLS) geschätzt:
	
	\begin{equation}
		\boldsymbol{\theta}_i = \arg\min_{\boldsymbol{\theta}} \sum_{k=1}^{N} \Phi_i(\mathbf{z}_k) \cdot \left(y_k - \boldsymbol{\theta}^T \tilde{\mathbf{x}}_k\right)^2
	\end{equation}
	
	Dies entspricht dem Standard-Least-Squares-Problem mit gewichteter Datenmatrix:
	
	\begin{equation}
		\boldsymbol{\theta}_i = \left(\tilde{\mathbf{X}}^T \mathbf{W}_i \tilde{\mathbf{X}}\right)^{-1} \tilde{\mathbf{X}}^T \mathbf{W}_i \mathbf{y}
	\end{equation}
	
	wobei $\mathbf{W}_i = \mathrm{diag}(\Phi_i(\mathbf{z}_1), \ldots, \Phi_i(\mathbf{z}_N))$.
	
	In der Implementierung wird \texttt{np.linalg.lstsq} auf die gewichtete Matrix $\mathbf{W}_i^{1/2} \tilde{\mathbf{X}}$ angewendet (numerisch stabiler).
	
	% ----------------------------------------
	\subsection{Lernalgorithmus -- Baumwachstum}
	% ----------------------------------------
	
	Der HILOMOT-Algorithmus wächst den Baum iterativ:
	
	\begin{lstlisting}[caption={HILOMOT-Lernalgorithmus (Pseudocode)},language={}
		]
		1. Starte mit einem einzigen Blatt (globales lineares Modell)
		2. Wiederhole bis max_leafs erreicht:
		a. Bestimme das "schlechteste" Blatt
		(hoechster gewichteter Trainingsfehler)
		b. Generiere Split-Kandidaten fuer dieses Blatt
		c. Evaluiere jeden Kandidaten:
		- Schaetze kappa
		- Berechne phi_left und phi_right
		- Passe lokale Modelle per WLS an
		- Berechne resultierenden Gesamtfehler
		d. Fuehre den besten Split durch (falls verbessernd)
		e. Falls kein Split den Fehler verbessert: Abbruch
	\end{lstlisting}
	
	Das schlechteste Blatt wird bestimmt durch:
	
	\begin{equation}
		i^* = \arg\max_i \sum_{k=1}^{N} \Phi_i(\mathbf{z}_k) \cdot \left(y_k - \hat{y}_i(\mathbf{x}_k)\right)^2
	\end{equation}
	
	% ----------------------------------------
	\subsection{Detaillierter Partitionierungsprozess}
	% ----------------------------------------
	
	\subsubsection*{1.~Bestimmung der aktiven Region}
	
	Nach jeder Iteration werden für alle Trainingspunkte die Blatt-Gültigkeiten berechnet:
	
	\begin{equation}
		\boldsymbol{\Phi}_i = [\Phi_i(\mathbf{z}_1), \ldots, \Phi_i(\mathbf{z}_N)]^T
	\end{equation}
	
	Dieser Vektor wirkt als \textbf{weiche Zuordnung} der Datenpunkte zu diesem Blatt.
	
	\subsubsection*{2.~Wahl des zu teilenden Blattes}
	
	Das Blatt mit dem grö\ss{}ten gewichteten Fehler wird ausgewählt. Ein Blatt mit vielen Punkten aber kleinem Fehler wird eher nicht gesplittet; ein Blatt mit starker Nichtlinearität wird bevorzugt gesplittet.
	
	\subsubsection*{3.~Bestimmung des Partitionszentrums}
	
	Das gewichtete Zentrum der Region wird berechnet:
	
	\begin{equation}
		\mathbf{c} = \frac{\sum_{k=1}^{N} \Phi_{\mathrm{worst}}(\mathbf{z}_k) \cdot \mathbf{z}_k}{\sum_{k=1}^{N} \Phi_{\mathrm{worst}}(\mathbf{z}_k)}
	\end{equation}
	
	Der Bias wird so gesetzt, dass die Hyperebene durch dieses Zentrum läuft:
	
	\begin{equation}
		b = -\mathbf{w}^T \mathbf{c}
	\end{equation}
	
	\subsubsection*{4.~Erzeugung von Split-Kandidaten}
	
	Für das zu teilende Blatt werden mögliche Richtungen erzeugt:
	\begin{enumerate}
		\item Alle Achsrichtungen $\mathbf{e}_1, \ldots, \mathbf{e}_n$
		\item Falls \texttt{oblique=True}: zusätzlich die Hauptachsenrichtung der gewichteten Kovarianzmatrix
	\end{enumerate}
	
	Die Kovarianzmatrix wird gewichtet aufgebaut:
	
	\begin{equation}
		\mathbf{C} = \sum_{k=1}^{N} \Phi_{\mathrm{worst}}(\mathbf{z}_k) (\mathbf{z}_k - \bar{\mathbf{z}})(\mathbf{z}_k - \bar{\mathbf{z}})^T
	\end{equation}
	
	\subsubsection*{5.~Richtungsoptimierung im 2D-Fall}
	
	Für zweidimensionale Eingänge wird die Richtung kontinuierlich optimiert:
	
	\begin{equation}
		\mathbf{w}(\theta) = [\cos\theta, \sin\theta]^T, \qquad
		J(\theta) = \mathrm{MSE}(\mathbf{y},\, \hat{\mathbf{y}}_{\mathrm{candidate}}(\theta))
	\end{equation}
	
	Minimiert via \texttt{scipy.optimize.minimize\_scalar} auf $[0, 2\pi]$.
	
	\subsubsection*{6.~Mindestdatenmasse}
	
	Nur wenn beide Kinder genügend effektive Datenmasse besitzen, wird der Kandidat weiter betrachtet:
	
	\begin{equation}
		\sum_k \Phi_{\mathrm{left}}(\mathbf{z}_k) \geq n_{\min}, \qquad
		\sum_k \Phi_{\mathrm{right}}(\mathbf{z}_k) \geq n_{\min}
	\end{equation}
	
	\subsubsection*{7.~Bildung der Kind-Gültigkeiten}
	
	\begin{equation}
		\Phi_{\mathrm{left}}(\mathbf{z}) = \Phi_{\mathrm{parent}}(\mathbf{z}) \cdot \Psi(\mathbf{z})
	\end{equation}
	\begin{equation}
		\Phi_{\mathrm{right}}(\mathbf{z}) = \Phi_{\mathrm{parent}}(\mathbf{z}) \cdot \bigl(1 - \Psi(\mathbf{z})\bigr)
	\end{equation}
	
	\subsubsection*{8.~Refit der lokalen Modelle}
	
	Für beide neuen Kindknoten werden sofort neue lokale lineare Modelle per WLS geschätzt:
	
	\begin{equation}
		\boldsymbol{\theta}_{\mathrm{left}} = \arg\min_{\boldsymbol{\theta}} \sum_k \Phi_{\mathrm{left}}(\mathbf{z}_k) \left(y_k - \boldsymbol{\theta}^T \tilde{\mathbf{x}}_k\right)^2
	\end{equation}
	
	\subsubsection*{9.~Bewertung des Kandidaten-Splits}
	
	\begin{equation}
		\hat{\mathbf{y}}_{\mathrm{new}} = \hat{\mathbf{y}}_{\mathrm{old}}
		- \Phi_{\mathrm{parent}} \hat{\mathbf{y}}_{\mathrm{parent}}
		+ \Phi_{\mathrm{left}} \hat{\mathbf{y}}_{\mathrm{left}}
		+ \Phi_{\mathrm{right}} \hat{\mathbf{y}}_{\mathrm{right}}
	\end{equation}
	
	Nicht der lokale, sondern der \textbf{globale Trainingsfehler} entscheidet, ob ein Split gut ist.
	
	\begin{figure}[H]
		\centering
		\includegraphics[width=0.85\textwidth]{hilomot_3d_partition_visualization.png}
		\caption{Visualisierung der HILOMOT-Partitionierung. Farbflächen zeigen dominante lokale Modelle, schwarze Linien die im Baum gespeicherten Trennhyperebenen.}
		\label{fig:partition}
	\end{figure}
	
	% ----------------------------------------
	\subsection{Modellfit und Vorhersage}
	% ----------------------------------------
	
	\subsubsection*{1.~Skalierung von Ein- und Ausgängen}
	
	\begin{equation}
		\mathbf{x}_s = \frac{\mathbf{x} - \mathbf{x}_{\min}}{\mathbf{x}_{\max} - \mathbf{x}_{\min}}, \qquad
		y_s = \frac{y - y_{\min}}{y_{\max} - y_{\min}}
	\end{equation}
	
	\subsubsection*{2.~Startmodell}
	
	Zu Beginn existiert nur ein einziges Blatt mit einem globalen linearen Modell:
	
	\begin{equation}
		\hat{y}^{(0)}(\mathbf{x}) = \boldsymbol{\theta}_{\mathrm{root}}^T \tilde{\mathbf{x}}
	\end{equation}
	
	\subsubsection*{3.~Iterativer Refit}
	
	Jeder akzeptierte Split ersetzt genau \textbf{ein} Blatt durch \textbf{zwei} neue Blätter. Der Algorithmus ist damit \textbf{greedy} und lokal inkrementell.
	
	\subsubsection*{4.~Vorhersage im fertigen Modell}
	
	Gültigkeiten werden normiert:
	
	\begin{equation}
		V_i(\mathbf{z}) = \frac{\Phi_i(\mathbf{z})}{\displaystyle\sum_{m=1}^{M} \Phi_m(\mathbf{z})}
	\end{equation}
	
	Die finale Vorhersage lautet:
	
	\begin{equation}
		\hat{y}(\mathbf{x}) = \sum_{i=1}^{M} V_i(\mathbf{z}) \cdot \hat{y}_i(\mathbf{x})
	\end{equation}
	
	\begin{figure}[H]
		\centering
		\includegraphics[width=0.85\textwidth]{hilomot_3d_visualization.png}
		\caption{Modellierte Antwortfläche des trainierten HILOMOT-Netzes über dem Eingangsraum.}
		\label{fig:surface}
	\end{figure}
	
	% ----------------------------------------
	\subsection{Wahl der besten Struktur}
	% ----------------------------------------
	
	\subsubsection*{Lokal beste Aufteilung}
	
	In jeder Iteration wird genau derjenige Split-Kandidat übernommen, der den globalen Trainings-MSE am stärksten reduziert:
	
	\begin{equation}
		s^* = \arg\min_{s \in \mathcal{S}_{\mathrm{current}}} \mathrm{MSE}(\mathbf{y}, \hat{\mathbf{y}}^{(s)})
	\end{equation}
	
	\subsubsection*{Abbruchkriterien}
	
	Das Baumwachstum endet, wenn mindestens eines der folgenden Kriterien erfüllt ist:
	\begin{itemize}
		\item \texttt{max\_leafs} ist erreicht
		\item kein zulässiger Split-Kandidat existiert
		\item kein Split verbessert den aktuellen globalen Fehler
		\item ein Kandidat verletzt die Mindestmasse \texttt{min\_samples\_per\_leaf}
	\end{itemize}
	
	\subsubsection*{Strukturwahl im Tool}
	
	Im Repository existieren dafür zwei Strategien:
	\begin{itemize}
		\item \texttt{optimize\_generalization.py}: sucht Konfigurationen mit kleinem Abstand zwischen Trainings- und Validierungsfehler
		\item \texttt{optimize\_accuracy.py}: sucht Konfigurationen mit gutem Fit bei gleichzeitiger Begrenzung von Oszillationen
	\end{itemize}
	
	% ============================================================
	\section{Gegenüberstellung: LOLIMOT vs. HILOMOT}
	% ============================================================
	
	\begin{center}
		\begin{tabular}{p{3.5cm}p{4.5cm}p{4.5cm}}
			\toprule
			\textbf{Aspekt} & \textbf{LOLIMOT} & \textbf{HILOMOT} \\
			\midrule
			Schnitte & Achsenparallel & Beliebig orientiert \\
			Trennfläche & $x_j = \mathrm{const}$ & $\mathbf{w}^\top \mathbf{x} + b = 0$ \\
			Anzahl Modelle & Tendenziell mehr bei schrägen Grenzen & Oft weniger nötig \\
			Training & Schnell, einfach & Aufwendiger \\
			Robustheit & Sehr robust & Empfindlicher \\
			Interpretierbarkeit & Rechtecke leicht vorstellbar & Etwas abstrakter \\
			Gut geeignet für & Achsennahe Trennungen, Prototyping & Gekoppelte Eingänge, kompakte Modelle \\
			\bottomrule
		\end{tabular}
	\end{center}
	
	\begin{merkbox}
		\textbf{Beide Verfahren} erzeugen glatte, interpretierbare Modelle aus lokalen linearen Bausteinen. HILOMOT enthält LOLIMOT als Spezialfall -- die Wahl hängt vom konkreten Problem ab.
	\end{merkbox}
	
	% ============================================================
	\section{Mathematische Zusammenfassung}
	% ============================================================
	
	\begin{center}
		\begin{tabular}{ll}
			\toprule
			\textbf{Symbol} & \textbf{Bedeutung} \\
			\midrule
			$M$ & Anzahl lokaler Modelle (Blätter) \\
			$\Phi_i(\mathbf{z})$ & Gültigkeitsfunktion des $i$-ten Blattes \\
			$\hat{y}_i(\mathbf{x})$ & Ausgabe des $i$-ten lokalen linearen Modells \\
			$\boldsymbol{\theta}_i$ & Parametervektor des $i$-ten Modells \\
			$\mathbf{w}$, $b$ & Normalenvektor und Bias der Trennhyperebene \\
			$\kappa$ & Schärfeparameter (Steilheit der Sigmoid-Funktion) \\
			$\rho$ & Glättigkeitsparameter (\texttt{smoothness}) \\
			$\mu^*$ & Demanded membership value \\
			$N$ & Anzahl der Trainingsdaten \\
			$g(\mathbf{x})$ & Weiche Entscheidungsfunktion am Gate-Knoten \\
			$\mathbf{C}$ & Gewichtete Kovarianzmatrix der Blattregion \\
			$\theta$ & Winkelparameter für 2D-Richtungsoptimierung \\
			\bottomrule
		\end{tabular}
	\end{center}
	
	% ============================================================
	\section{Implementierung}
	% ============================================================
	
	\subsection{Projektstruktur}
	
	\begin{lstlisting}[caption={Verzeichnisstruktur des LMNTool-Projekts},language={}
		]
		LMNTool/
		|-- src/
		|   |-- main.py                    # Skript-Einstiegspunkt
		|   |-- HILO_GUI.py                # PyQt5-GUI-Hauptfenster
		|   |-- hilomot.py                 # Oeffentliche Hilomot-Klasse
		|   |-- hilomot_regressor.py       # Kernimplementierung
		|   |-- dataset.py                 # Datensatzverwaltung
		|   |-- globalmodel.py             # Globales Modell-Hilfsobjekt
		|   |-- parameter_export.py        # JSON/Text-Export
		|   |-- optimize_accuracy.py       # Hyperparameter-Optimierung
		|   |-- optimize_smoothness.py     # Hyperparameter-Optimierung
		|   |-- optimize_sharpness.py      # Hyperparameter-Optimierung
		|   +-- optimize_generalization.py # Kreuzvalidierungsbasiert
		|-- run_gui.py                     # Startskript GUI
		|-- run_main.py                    # Startskript Skriptmodus
		|-- requirements.txt               # Python-Abhaengigkeiten
		+-- README.md                      # Kurzbeschreibung
	\end{lstlisting}
	
	\subsection{Klassen und Module}
	
	\subsubsection*{\texttt{HilomotRegressor} (\texttt{hilomot\_regressor.py})}
	
	Kernklasse des Algorithmus.
	
	\begin{center}
		\begin{tabular}{p{5.5cm}p{7.5cm}}
			\toprule
			\textbf{Methode} & \textbf{Beschreibung} \\
			\midrule
			\texttt{fit(X, y)} & Trainiert das Modell auf den Daten \texttt{X} (N$\times$2) und Ausgabe \texttt{y} (N,) \\
			\texttt{predict(X)} & Berechnet die Modellvorhersage für neue Eingaben \\
			\texttt{score(X, y)} & Gibt den R\textsuperscript{2}-Wert auf dem Datensatz zurück \\
			\texttt{\_compute\_leaf\_validities(Z)} & Berechnet die Gültigkeitswerte $\Phi_i$ aller Blätter \\
			\texttt{\_choose\_best\_split(...)} & Evaluiert alle Split-Kandidaten \\
			\texttt{\_optimize\_split\_direction\_2d(...)} & Optimiert die Split-Richtung im 2D-Fall \\
			\texttt{\_estimate\_split\_kappa(...)} & Schätzt $\kappa$ für einen Kandidaten-Split \\
			\texttt{\_adjust\_kappa(...)} & Passt $\kappa$ basierend auf \texttt{demanded\_membership\_value} an \\
			\texttt{to\_dict()} / \texttt{from\_dict()} & Serialisierung als Python-Dictionary \\
			\texttt{save\_mat()} / \texttt{from\_mat()} & Speichern/Laden im MATLAB \texttt{.mat}-Format \\
			\bottomrule
		\end{tabular}
	\end{center}
	
	\subsubsection*{\texttt{HilomotTreeNode} (\texttt{hilomot\_regressor.py})}
	
	Knotenobjekt des Baums.
	
	\begin{center}
		\begin{tabular}{p{3.5cm}p{9cm}}
			\toprule
			\textbf{Attribut} & \textbf{Beschreibung} \\
			\midrule
			\texttt{w} & Gewichtsvektor $[\mathbf{w}; b]$ der Trennhyperebene \\
			\texttt{kappa} & Schärfeparameter $\kappa$ \\
			\texttt{center} & Zentrum der Partition \\
			\texttt{parameter} & Parametervektor $\boldsymbol{\theta}_i$ des lokalen Modells \\
			\texttt{is\_leaf} & \texttt{True} wenn Blattknoten \\
			\texttt{left} / \texttt{right} & Kindknoten \\
			\bottomrule
		\end{tabular}
	\end{center}
	
	% ============================================================
	\section{Tool-Beschreibung: LMNTool GUI}
	% ============================================================
	
	\subsection{Starten der Anwendung}
	
	\begin{lstlisting}[language=bash,caption={GUI starten}]
		python run_gui.py
	\end{lstlisting}
	
	Oder innerhalb der \texttt{src/}-Umgebung:
	
	\begin{lstlisting}[language=bash]
		python -m src.HILO_GUI
	\end{lstlisting}
	
	\subsection{Datenladen}
	
	Die GUI akzeptiert \textbf{CSV-Dateien} im 3-Spalten-Format:
	
	\begin{center}
		\begin{tabular}{lll}
			\toprule
			\textbf{Spalte} & \textbf{Physikalische Grö	tfamily{sse}} & \textbf{Beschreibung} \\
			\midrule
			Spalte~1 (\texttt{y}) & Tilt angle & Ausgangsgrö	tfamily{sse (Zielwert)} \\
			Spalte~2 (\texttt{x1}) & Shear force & Eingangsgrö	tfamily{sse~1} \\
			Spalte~3 (\texttt{x2}) & Normal force & Eingangsgrö	tfamily{sse~2} \\
			\bottomrule
		\end{tabular}
	\end{center}
	
	Das Trennzeichen wird automatisch erkannt (\texttt{;}, \texttt{,} oder Tab).
	
	\subsection{Modellparameter}
	
	\begin{center}
		\begin{tabular}{p{3.5cm}p{2cm}p{1.5cm}p{5cm}}
			\toprule
			\textbf{Parameter} & \textbf{GUI-Element} & \textbf{Standard} & \textbf{Beschreibung} \\
			\midrule
			Max Leaves & Spinbox & 15 & Maximale Anzahl lokaler Modelle \\
			Min Samples per Leaf & Spinbox & 6 & Minimale effektive Datenpunkte pro Blatt \\
			Smoothness ($\rho$) & Spinbox & 0{,}8 & überlappungsgrad der Regionen \\
			Membership Value ($\mu^*$) & Spinbox & 0{,}4 & Mindest-Gültigkeit eines Datenpunkts \\
			Oblique Splits & Checkbox & aktiv & Aktiviert schräge Trennhyperebenen \\
			\bottomrule
		\end{tabular}
	\end{center}
	
	\textbf{Empfehlungen:}
	\begin{itemize}
		\item Für komplexe nichtlineare Daten: höheres \texttt{max\_leafs} (10--20), niedrigeres \texttt{smoothness} (0{,}3--0{,}6)
		\item Für glatte Approximation: höheres \texttt{smoothness} (0{,}8--1{,}5), niedrigeres \texttt{max\_leafs}
		\item \texttt{oblique=True} liefert in der Regel bessere Ergebnisse bei schrägen Strukturen im Eingangsraum
	\end{itemize}
	
	\subsection{Training}
	
	\begin{itemize}
		\item \textbf{Train Model}: Trainiert das Modell im Hintergrund-Thread. Die GUI bleibt reaktionsfähig.
		\item \textbf{Run Full Analysis}: Vollständiges Training, erstellt alle Plots und exportiert das Modell automatisch.
	\end{itemize}
	
	Nach dem Training werden folgende Metriken angezeigt:
	
	\begin{center}
		\begin{tabular}{p{3cm}p{5cm}p{4.5cm}}
			\toprule
			\textbf{Metrik} & \textbf{Formel} & \textbf{Beschreibung} \\
			\midrule
			MSE & $\frac{1}{N}\sum(y_i - \hat{y}_i)^2$ & Mittlerer quadratischer Fehler \\
			MAE & $\frac{1}{N}\sum|y_i - \hat{y}_i|$ & Mittlerer absoluter Fehler \\
			Correlation & $r_{y,\hat{y}}$ & Pearson-Korrelationskoeffizient \\
			R\textsuperscript{2} & $1 - \frac{SS_{\mathrm{res}}}{SS_{\mathrm{tot}}}$ & Bestimmtheitsma\ss{} \\
			Local Models & $M$ & Anzahl der Blattknoten im Baum \\
			\bottomrule
		\end{tabular}
	\end{center}
	
	\subsection{Visualisierungen}
	
	\begin{center}
		\begin{tabular}{p{4cm}p{9cm}}
			\toprule
			\textbf{Plot-Typ} & \textbf{Beschreibung} \\
			\midrule
			3D Surface & Vorhersagefläche $\hat{y}(x_1, x_2)$ mit Trainingsdaten \\
			Training Data 3D & Nur die Rohdaten im 3D-Raum \\
			Partition Map & 2D-Partitionierungsplot mit Trennlinien \\
			Correlation Plot & Streudiagramm $y$ vs. $\hat{y}$ \\
			Shear Force Cuts & Schnitte durch das Modell bei verschiedenen Normalkraft-Werten \\
			\bottomrule
		\end{tabular}
	\end{center}
	
	\begin{figure}[H]
		\centering
		\includegraphics[width=0.85\textwidth]{hilomot_shear_cup_plots.png}
		\caption{Schnitte durch die modellierte Fläche entlang ausgewählter Scherkraft-Niveaus.}
		\label{fig:cuts}
	\end{figure}
	
	\subsection{Export}
	
	\begin{center}
		\begin{tabular}{ll}
			\toprule
			\textbf{Format} & \textbf{Inhalt} \\
			\midrule
			JSON & Vollständiger Baum: Knoten, Gewichte, Kappa-Werte, Modellparameter \\
			Text & Lesbare Zusammenfassung: Metriken, lokale Modellparameter \\
			\bottomrule
		\end{tabular}
	\end{center}
	
	% ============================================================
	\section{Skript-Nutzung (\texttt{main.py})}
	% ============================================================
	
	Das Skript \texttt{src/main.py} trainiert das Modell auf dem konfigurierten CSV-Datensatz und zeigt:
	\begin{enumerate}
		\item \textbf{3D-Vorhersagefläche} mit Trainingsdaten
		\item \textbf{Oblique-Partitionierungsplot} mit Trennlinien und LM-Beschriftungen
	\end{enumerate}
	
	\begin{lstlisting}[language=bash]
		python src/main.py
	\end{lstlisting}
	
	Konfigurierbare Konstanten in \texttt{main.py}:
	
	\begin{lstlisting}[language=python,caption={Konfiguration in main.py}]
		csv_path = r"\\server\pfad\zu\datei.csv"   # Datenpfad
		max_leafs = 10                              # Max. Blaetter
		min_samples_per_leaf = 10                  # Min. Samples/Blatt
		oblique = True                             # HILOMOT (True) oder LOLIMOT (False)
	\end{lstlisting}
	
	% ============================================================
	\section{Parameter-Referenz}
	% ============================================================
	
	\begin{center}
		\begin{tabular}{p{4cm}p{1.2cm}p{2.2cm}p{2cm}p{4.5cm}}
			\toprule
			\textbf{Parameter} & \textbf{Typ} & \textbf{Bereich} & \textbf{Empfehlung} & \textbf{Effekt} \\
			\midrule
			\texttt{max\_leafs} & int & 1--100 & 8--20 & Mehr Blätter $\Rightarrow$ höhere Flexibilität, Overfitting-Risiko \\
			\texttt{min\_samples\_per\_leaf} & int & 1--50 & 5--15 & Zu klein $\Rightarrow$ instabile Modelle \\
			\texttt{smoothness} & float & 0{,}1--2{,}0 & 0{,}4--1{,}0 & Kleiner $\Rightarrow$ schärfere übergänge \\
			\texttt{demanded\_membership\_value} & float & 0{,}1--1{,}0 & 0{,}4--0{,}8 & Höher $\Rightarrow$ stärkere Lokalisierung \\
			\texttt{oblique} & bool & -- & \texttt{True} & \texttt{True}: HILOMOT; \texttt{False}: LOLIMOT \\
			\bottomrule
		\end{tabular}
	\end{center}
	
	% ============================================================
	\section{Literatur}
	% ============================================================
	
	\begin{enumerate}
		\item \textbf{Nelles, O.} (2001). \textit{Nonlinear System Identification}. Springer, Berlin.\\
		Grundlagenwerk zu lokalen Modellnetzen, LOLIMOT und HILOMOT.
		
		\item \textbf{Nelles, O., Fink, A., Isermann, R.} (2000). \textit{Local Linear Model Trees (LOLIMOT) Toolbox for Nonlinear System Identification}. Proc. IFAC Symposium on System Identification (SYSID).
		
		\item \textbf{Hartmann, B., Nelles, O.} (2009). \textit{Hierarchical Local Model Trees for Design of Experiments in the Framework of Global Model Networks}. Proc. 15th IFAC Symposium on System Identification.
		
		\item \textbf{LMNTool MATLAB Toolbox v1.5.2.1} -- Universität Siegen, Lehrstuhl für Regelungs- und Systemtechnik.\\
		Referenzimplementierung, auf der diese Python-Portierung basiert.
	\end{enumerate}
	
\end{document}